% ===================================================================
%  TABELL Nr. 2 — De mënschleche Kierper. — D'Paus. — D'Spiller.
%  Traduction 2026 d'apres texto/io, controlee sur texto/fr, texto/de
%  et texto/nl. Consigne : voir texto/lb/10-tabell-01.tex.
%
%  LE TABLEAU DU CORPS EST LE BANC D'ESSAI DE CETTE COLONNE, et il
%  fait pour le luxembourgeois ce qu'aucun tableau n'avait pu faire
%  pour les quatorze precedentes : IL MONTRE LES DEUX MOITIES DE LA
%  LANGUE SUR LA MEME PAGE.
%
%  La premiere moitie, le corps, est germanique de bout en bout et
%  n'a pas un mot francais : « Kapp », « Stir », « Aa », « Nues »,
%  « Baken », « Oueren », « Mond », « Kënn », « Zong », « Hals »,
%  « Broscht », « Häerz », « Longen », « Rippen », « Réck »,
%  « Bauch », « Mo », « Knéi », « Wued », « Knéchel », « Fouss »,
%  « Sool », « Fersen ». Quarante noms, zero emprunt roman.
%
%  La seconde moitie, les jeux, est francaise partout ou le jeu s'est
%  appris a l'ecole ou au club : « Kroquet », « Bilboquet »,
%  « Raquette », « Volant », « Trapez », « Tennis », « Foussball »,
%  « Vëlo », « boxen », « fechten », « Fleuret », « Hantelen ». Et
%  reste luxembourgeois tout ce qui s'apprend dans la cour :
%  « Verstoppes » cache-cache, « Sprangsäel » la corde a sauter,
%  « Stëlzen » les echasses, « Murmelen » les billes, « Kräisel » la
%  toupie, « Blannemann » le colin-maillard, « Fräschespill » le jeu
%  de la grenouille.
%
%  LA MESURE GALICIENNE DOIT DONC ETRE DOUBLEE ICI. Le galicien
%  mesurait la distance a UNE langue voisine, et la trouvait maximale
%  au corps, minimale a l'ecole. Le luxembourgeois a DEUX voisins, et
%  les deux distances varient EN SENS CONTRAIRE : la ou il s'eloigne
%  du francais il se rapproche de l'allemand, et inversement. Une
%  seule page suffit a le lire, parce qu'elle contient les deux
%  extremes.
%
%  ET LA NOTE DU FEUILLET 17 TOMBE JUSTE ICI. L'ido y dit que les
%  jeux d'enfants portent des noms « pure nacionala » et qu'il se
%  reserve soit de nommer par l'action, soit de prendre le nom
%  national. Le luxembourgeois fait exactement cela, mais sans
%  choisir : il prend le nom francais pour le jeu venu de France et
%  nomme par l'action celui de la cour d'ecole.
% ===================================================================

%%K t02-apar-1 apar
\VUsaut{4.0mm}
\VUtitre{15.0pt}{60}{TABELL Nr. 2}
\VUsaut{2.0mm}
\VUtitre{11.4pt}{0}{De mënschleche Kierper. --- D'Paus.}
\VUtitre{11.4pt}{0}{D'Spiller.}

%%K t02-tit-0 sub
\VUtitre{13.2pt}{0}{De mënschleche Kierper.}
\VUsaut{2.4mm}
\VUcentre{10.2pt}{0}{\textit{(Eng einfach Naturkonschtlektioun.)}}
\VUsaut{3.0mm}

%%K t02-01-1 p
1. --- D'Haaptdeeler vum mënschleche Kierper sinn: de \VUgras{Kapp,} de
\VUgras{Rimpf} an d'\VUgras{Glidder.}

%%K t02-tit-1 sub
\VUpk{10.2pt}{40}{De Kapp.}
\VUsaut{3.0mm}

%%K t02-02-1 p
2. --- De Kapp huet zwee Deeler: de \VUgras{Schädel,} an deem d'\VUgras{Gehir} ass an
deen d'\VUgras{Hoer} bedecken, an d'\VUgras{Gesiicht.}

%%K t02-03-1 p
3. --- Op eiser Zeechnung gesi mer weder de Schädel nach d'Gehir; mä mir gesi ganz
kloer déi wichteg Deeler vum Gesiicht.

%%K t02-04-1 p
4. --- Do ass fir d'éischt d'\VUgras{Stir,} déi e mächtege Verstand uweist, en Denken
dat ëmmer schafft. D'\VUgras{Schléifen} si ganz fräi. D'\VUgras{Aa} ass liewweg a
wäit op. Eng dicht \VUgras{Aebra} a laang \VUgras{Wimperen} schützen et.

%%K t02-05-1 p
5. --- Do sinn nach d'\VUgras{Nues} an d'\VUgras{Nueslächer.} D'Nues déngt eis fir ze
richen a fir ze otmen. Op béide Säite vun der Nues sinn d'\VUgras{Baken,} méi wäit
ewech d'\VUgras{Oueren.} Den ënneschten Deel vum Gesiicht besteet aus dem
\VUgras{Mond} an dem \VUgras{Kënn.}

%%K t02-06-1 p
6. --- Mir gesi se net ganz gutt, well de \VUgras{Schnauzert} an de
\VUgras{Bakenbaart} se eis verstoppen. Mä mir wëssen, datt de Mond déi zwou
\VUgras{Lëpsen} huet, den \VUgras{Uewerkifer} an den \VUgras{Ënnerkifer} mat hire
\VUgras{Zänn,} an d'\VUgras{Zong.} Mam Mond schwätze mir an iesse mir. Mat der Zong
schmaache mer d'Iessen.

%%K t02-07-1 p
7. --- D'Aa, d'Ouer, d'Nues an de Mond sinn, sou wéi d'Fanger, d'Organer vun de
fënnef Sënner: d'\VUgras{Gesiicht,} d'\VUgras{Gehéier,} de \VUgras{Geroch,} de
\VUgras{Geschmaach,} den \VUgras{Tastsënn.}

%%K t02-08-1 p
8. --- De Kapp ass also ee vun den interessantsten Deeler vum mënschleche Kierper.

%%K t02-tit-2 sub
\VUsaut{4.4mm}
\VUpk{10.2pt}{40}{De Rimpf.}
\VUsaut{3.0mm}

%%K t02-09-1 p
9. --- Den \VUgras{Hals} verbënnt de Kapp mam Rimpf. Hien huet d'\VUgras{Kiel} an den
\VUgras{Nak.}

%%K t02-10-1 p
10. --- De Rimpf huet och vill wichteg Organer. An der \VUgras{Broscht} sinn
d'\VUgras{Longen,} mat deene mir otmen, an d'\VUgras{Häerz,} dat de Mëttelpunkt vum
Liewen ass. Wann d'Häerz ophält ze schloen, stierft d'liewegt Wiesen.

%%K t02-11-1 p
11. --- D'\VUgras{Rippen} halen d'Broscht.

%%K t02-12-1 p
12. --- Déi aner Deeler vum Rimpf sinn d'\VUgras{Säit,} de \VUgras{Réck,}
d'\VUgras{Schëlleren} an de \VUgras{Bauch.} Fir ze schlofe leeë mir eis op d'Säit
oder op de Réck; d'Laaschte droe mir op der Schëller.

%%K t02-13-1 p
13. --- Am Bauch sinn de \VUgras{Mo} an d'\VUgras{Gedärm.} Si hëllefen eis, d'Iessen
ze verdauen.

%%K t02-tit-3 sub
\VUsaut{4.4mm}
\VUpk{10.2pt}{40}{D'Glidder.}
\VUsaut{3.0mm}

%%K t02-14-1 p
14. --- Mir hu véier \VUgras{Glidder}: zwee \VUgras{Äerm} an zwee \VUgras{Been.} Mat
den Äerm huele mir, halen, hiewen a sammele mir d'Saachen; mat de Beng stinn, ginn a
lafe mir.

%%K t02-15-1 p
15. --- D'\VUgras{Schëller} verbënnt den Aarm mam Kierper. E ganz staarke
\VUgras{Muskel,} de Bizeps, beweegt den Aarm, deen den \VUgras{Ielebou} mam
\VUgras{Ënneraarm} verbënnt. D'\VUgras{Handgelenk} bilt d'Gelenk vun der
\VUgras{Hand,} déi mat de \VUgras{Fanger} an den \VUgras{Neel} ophält. Op der Hand
gesi mer eng \VUgras{Vene} (an eng Arterie), an där d'\VUgras{Blutt} leeft.

%%K t02-16-1 p
16. --- D'Deeler vum Been sinn: den \VUgras{Uewerschenkel,} d'\VUgras{Knéi} an de
\VUgras{Knéikeel,} d'\VUgras{Wued,} vun där d'\VUgras{Fleesch} vun der \VUgras{Haut}
bedeckt ass, de \VUgras{Knéchel} an de \VUgras{Fouss.} D'\VUgras{Schanken} vum Been
si ganz déck a ganz staark. Um Enn vum Fouss sinn d'\VUgras{Zeeën.} Déi aner Deeler
sinn d'\VUgras{Sool} an d'\VUgras{Fersen.}

%%K t02-17-1 p
17. --- Dat ass déi bal komplett Beschreiwung vum mënschleche Kierper.

%%K t02-17-2 p
\textit{Bemierkung.} --- \textit{Mat Hëllef vum Ausdrock: „ech hu Wéi um... (um Kapp
asw.)“ kann de Léierer dëse ganze Wuertschatz nach eng Kéier benotzen, ënnert dem
Gesiichtspunkt vum gudden oder schlechte} \VUgras{Kierperzoustand.}

%%K t02-tit-4 sub
\VUsaut{5.0mm}
\VUpk{10.2pt}{40}{Gymnastik.}
\VUsaut{2.4mm}
\VUcentre{10.2pt}{0}{\textit{(Erzielung vun engem vun de Schüler.)}}
\VUsaut{3.0mm}

%%K t02-18-1 p
18. --- Fir gelenkeg, flénk a gesond ze sinn, musse mir eis Been an eis Äerm üben.
Dofir spille mir a maache mir \VUgras{Gymnastik.}

%%K t02-19-1 p
19. --- Während der Woch fannen dës zwou Übungen am Haff vum Lycée statt (kuck
d'Tabell Nr.~1). Mä donneschdes a sonndes gi mer op eng grouss Wiss, wou mer Plaz hu
fir ze lafen an eis ze ameséieren.

%%K t02-20-1 p
20. --- Eis Meeschteren an eis Elteren begleeden eis heiansdo dohin; mä et ass de
\VUgras{Gymnastiksléierer} \textsuperscript{(12)}, deen d'Übunge leet.

%%K t02-21-1 p
21. --- Op der Wiss, lénks vum Agank, stinn d'\VUgras{Turngeräter}
\textsuperscript{(1)}: mir klamme mat der \VUgras{Leeder} \textsuperscript{(2)},
maachen d'Kopp-ënnescht un de \VUgras{Ringen} \textsuperscript{(3)} an um
\VUgras{Trapez} \textsuperscript{(4)}, mir zéien eis mat den Hänn bis uewen op
d'\VUgras{Strackleeder} \textsuperscript{(5)} oder op d'\VUgras{Knuetesäel}
\textsuperscript{(6)}.

%%K t02-22-1 p
22. --- Mëttlerweil sprangen eis Kameraden iwwert d'\VUgras{Holzpäerd}
\textsuperscript{(7)} oder iwwert d'\VUgras{Barren} \textsuperscript{(11)}. Anerer
\VUgras{boxen} \textsuperscript{(8)} oder \VUgras{fechten} \textsuperscript{(9)} mat
Mask a mat der \VUgras{Fleuret.} Anerer schlussendlech \VUgras{hiewen Hantelen}
\textsuperscript{(10)}.

%%K t02-tit-5 sub
\VUsaut{4.6mm}
\VUpk{10.2pt}{40}{D'Spiller.}
\VUsaut{3.0mm}

%%K t02-23-1 p
23. --- Ronderëm d'Turner spillen d'Bouwen an d'Meedercher verschidde
\VUgras{Spiller.} D'Meedercher ameséiere sech mat hirem \VUgras{Rëf}
\textsuperscript{(14)}; si sprange mam \VUgras{Sprangsäel} \textsuperscript{(13)} oder
invitéieren hir Bridder a Cousinen op eng \VUgras{Kroquet}-Partie
\textsuperscript{(15)}. Si fanne Freed drun, dem Géigner seng \VUgras{Kugel} mat hirem
\VUgras{Holzhummer} ewechzestoussen, grad an deem Moment, wou en denkt, hie kéim
liicht duerch den Archbou!

%%K t02-24-1 p
24. --- D'Bouwen huele léiwer méi muskulös Spiller. Si spille gär \VUgras{Keelen}
\textsuperscript{(16)}, déi si geschéckt mat enger \VUgras{Kugel} \textsuperscript{(17)}
ëmwerfen. Si veruechten och de \VUgras{Kräisel} \textsuperscript{(18)} an de
\VUgras{Bilboquet} \textsuperscript{(19)} net, oder esouguer d'\VUgras{Fräschespill}
\textsuperscript{(20)}. Mä hir léifste Spiller sinn d'\VUgras{Murmelen}
\textsuperscript{(21)} an de \VUgras{Mauerball} \textsuperscript{(22)}, well d'Mauer
vun der \VUgras{Serre} \textsuperscript{(26)} ganz gutt gëeegent ass, fir de liichte
Ball zréckspréngen ze loossen.

%%K t02-25-1 p
25. --- De klenge Jang, deen op senge \VUgras{Stëlzen} \textsuperscript{(24)} geet,
ass ganz geschéckt a kann d'Gläichgewiicht ganz gutt halen. Nieft him spillen e puer
Kameraden \VUgras{Verstoppes} \textsuperscript{(25)} an där Serre an hannert der
\VUgras{Heck.} Anerer huelen léiwer d'\VUgras{Waarmhand} \textsuperscript{(28)} oder de
\VUgras{Volant} \textsuperscript{(29)}, deen liicht vun enger \VUgras{Raquette} bis
bei déi aner flitt.

%%K t02-noto-1 noto
\VUnotes{143.2mm}{%
(*) D'Kannerspiller hu ganz dacks reng national Nimm. Mir hunn eis Méi ginn, si esou
gutt wéi méiglech op Ido ze nennen, entweder andeems mir eis vun der Handlung
inspiréiere loossen, déi si kennzeechent, oder andeems mir den nationalen Numm an
dësem oder an deem Land wielen, wann en net ze ongerecht ass. Zum Beispill:
d'\VUgras{Faassspill} (däitsch a franséisch) ass d'\VUgras{Fräschespill}
\textsuperscript{(20)} (englesch), an deem d'Spiller probéieren, hir ronn Scheiwen
duerch de Mond vun engem Metallfräsch ze werfen, dee mat oppenem Maul op enger
duerchlächerter Këscht (engem Faass) sëtzt.
De \VUgras{Schëllersprong} \textsuperscript{(30)} ënnerscheet sech vum
\VUgras{Récksprong} nëmmen doduerch: fir iwwert de Kapp ze sprangen, stäipen d'Spiller
d'Hänn op d'Schëllere vun hirem Partner, wärend si beim „\VUgras{Récksprong}“
d'Hänn nëmmen op säi Réck stäipen. --- Oder e puer vun den Deelhueler béie sech, an
déi aner sprangen iwwert hire Réck, andeems si d'Hänn op d'Schëller stäipen; déi
éischt musse de Stouss aushalen ouni anzeknécken, déi zweet musse sprangen ouni
d'Gläichgewiicht ze verléieren (kuck d'Tabell selwer).%
}

%%K t02-26-1 p
26. --- An der Mëtt vun der Wiss stinn zwee Beem. Um Fouss vun engem vun deenen zwee
hu siwen oder aacht Bouwen eng Partie \VUgras{Schëllersprong} \textsuperscript{(30)}
(*) organiséiert. Dat Spill ass net an alle Länner bekannt; mä jidderee weess, wat de
\VUgras{Récksprong} ass, mat deem d'Kanner dës Kéier net spillen.

%%K t02-tit-6 sub
\VUsaut{5.0mm}
\VUpk{10.2pt}{40}{D'Spiller an der fräier Loft.}
\VUsaut{3.4mm}

%%K t02-27-1 p
27. De \VUgras{Blannemann} \textsuperscript{(31)} ass och ganz ameséiereg;
Meedercher a Bouwe leeën ofwiesselnd d'\VUgras{Aendëchel} un an huele déi
Ongeschéckter, déi sech fänke loossen.

%%K t02-28-1 p
28. --- An der Mëtt vun der Wiss, do ass e ganz franséischt mä e bëssen hefteg
Spill: dat ass d'\VUgras{Bierespill} \textsuperscript{(32)}, vill méi haart wéi dat
esou roueg Spill vum \VUgras{Draach} \textsuperscript{(33)} oder esouguer wéi dat
englescht Spill \VUgras{Tennis} \textsuperscript{(34)}.

%%K t02-29-1 p
29. --- Dee leschte Sport ass d'Freed vun de groussen Schüler. Eis Léierer selwer
spillen en gär. Hie verlaangt vill Geschécklechkeet a Flinkheet, an hie gefält mer
ganz gutt. D'Spiller mat der \VUgras{Schaukel} \textsuperscript{(35)} loossen ech de
Meedercher.

%%K t02-30-1 p
30. --- En anert englescht Spill hunn ech net gär, well et vläicht geféierlech ginn
kéint.

%%K t02-30-1b p
Ech mengen de \VUgras{Foussball} \textsuperscript{(36)}, deen menge grousse Brudder
frou mécht. Ech hunn eist aalt \VUgras{Barrespill} \textsuperscript{(38)} méi gär, dat
grad esou gesond a wierklech manner brutal ass.

%%K t02-tit-7 sub
\VUsaut{4.6mm}
\VUpk{10.2pt}{40}{Vëlofueren a Ballonspill.}
\VUsaut{3.0mm}

%%K t02-31-1 p
31. --- Ech fueren och gär mam \VUgras{Vëlo} \textsuperscript{(37)} ronderëm d'Wiss.

%%K t02-32-1 p
32. --- Nächst Joer léieren ech d'\VUgras{Reiden} \textsuperscript{(39)}. Wéi schéin
ass e Päerd, dat zierlech schrëtt oder trëppelt, an dat am Galopp iwwert
d'Hindernisser sprangt!

%%K t02-33-1 p
33. --- Am Summer léiere mir d'\VUgras{Schwammen} \textsuperscript{(40)}. Mir
tauchen a bueden eis mat Freed am kloren a frëschen Waasser vum Floss. Heiansdo
loosse mer schlussendlech e pabeierne \VUgras{Ballon} \textsuperscript{(41)}
opsteigen, deen majestéitesch an d'Loft klëmmt, esou bal wéi en mat waarmer Loft
gefëllt ass.

%%K t02-34-1 p
34. --- Et gëtt nach vill aner Spiller, mä ech géif net fäerdeg ginn, se opzezielen,
an een erkennt aus dëser Zesummefaassung, datt ee sech donneschdes op eiser schéiner
Wiss net vill langweilt.

%%K t02-34-2 p
N.-B. --- \textit{De Léierer wäert dëst Kapitel liicht ergänzen, andeem en all eenzelt
vun de wichtegste Spiller méi ausféierlech beschreift.}
